\documentclass[preprint,12pt]{elsarticle}
\usepackage{amsmath}
\usepackage{booktabs}
\usepackage{graphicx}
\usepackage{hyperref}
\usepackage{float}
\usepackage{xcolor}
\usepackage{makecell}
\usepackage{caption}
\usepackage{longtable}
\usepackage{booktabs}

\captionsetup[table]{font=normal}
\biboptions{authoryear}

\setlength{\parskip}{0.5em}
\renewcommand{\arraystretch}{1.2}

\journal{Economics Letters}

\begin{document}
\begin{center}
    {\Large\bfseries Supplementary Material}\\[0.3em]
    {\normalsize\itshape Weakening Environmental Enforcement and Homicides: Evidence from Indigenous Territories in Brazil}
\end{center}

\vspace{1em}

\section{Data} \label{ap:data}

The paper is based on three constructed variables: homicide rate per 100,000 inhabitants, a binary variable for gold mining reserves in the municipality, and a binary variable for indigenous territories in the municipality. Therefore, it is worth highlighting their composition and possible limitations.     

\subsection*{A.1. Homicide Rate}

The homicide rate is the simple ratio of homicides to population, multiplied by 100,000.

\begin{itemize}
    \item In the numerator, homicide is defined as the sum of three ICD-10 cause-of-death categories, by municipality of occurrence: assaults (X85–Y09), events of undetermined intent (Y10–Y34), and legal interventions (Y35). The raw data was sourced from SIM/DATASUS.
    
    \item In the denominator, population figures are estimated by the Brazilian Institute of Geography and Statistics (IBGE) for the Federal Court of Accounts (TCU), within the scope of the Municipalities Participation Fund (FPM), by municipality of residence. The data is available on the IBGE website.
\end{itemize}

It is worth mentioning some points. First, the numerator is measured by the municipality of \textit{occurrence} while the denominator is measured by the municipality of \textit{residence}. Although SIM/DATASUS allows, by default, the retrieval of deaths by municipality of residence, we opted for a perspective based on occurrence. The underlying reasoning is that municipalities where Indigenous territories overlap gold reserves often have a substantial share of inhabitants whose residence is outside the municipality –- that is, low-skilled workers who migrate from their places of residence in search of increased income from high rents of natural resources. In contrast, only IBGE's population estimates are available at the municipal level, and they only consider the perspective of residence.

Second, other ICD-10 cause-of-death categories could be used to build the homicide variable -- or even some category could be removed. The existing literature often addresses homicides as a combination of one or more of the mentioned categories, in addition to intentional self-harm (X60-X84). We opted to include the three most used ICD-10 cause-of-death categories. A comprehensive table with all ICD-10 definitions ranging from X60 to Y35 is displayed for consultation.

\begingroup
\small
\renewcommand{\arraystretch}{1.2}
\setlength{\tabcolsep}{3.5pt}
% icd10_table.tex
% ICD-10 Codes X60-Y35 grouped by official WHO blocks
% Required packages in main document:
%   \usepackage{longtable}
%   \usepackage{booktabs}

\begin{longtable}{p{2.2cm} p{10.3cm}}

\caption{International Statistical Classification of Diseases and Related Health Problems - 10th Revision (ICD-10): Codes X60--Y35}
\label{tab:icd10_x60_y35} \\

\toprule
\textbf{Code} & \textbf{Description} \\
\midrule
\endfirsthead

\multicolumn{2}{l}{\small\itshape (Continued from previous page)} \\[2pt]
\toprule
\textbf{Code} & \textbf{Description} \\
\midrule
\endhead

\midrule
\multicolumn{2}{r}{\small\itshape (Continued on next page)} \\
\endfoot

\bottomrule
\endlastfoot

% ================================================================
% BLOCK X60-X84  Intentional self-harm
% ================================================================
\multicolumn{2}{l}{\textbf{X60--X84 \quad Intentional Self-Harm}} \\
\midrule
\texttt{X60} & Intentional self-poisoning by and exposure to nonopioid analgesics, antipyretics and antirheumatics \\
\texttt{X61} & Intentional self-poisoning by and exposure to antiepileptic, sedative-hypnotic, antiparkinsonism and psychotropic drugs, not elsewhere classified \\
\texttt{X62} & Intentional self-poisoning by and exposure to narcotics and psychodysleptics (hallucinogens), not elsewhere classified \\
\texttt{X63} & Intentional self-poisoning by and exposure to other drugs acting on the autonomic nervous system \\
\texttt{X64} & Intentional self-poisoning by and exposure to other and unspecified drugs, medicaments and biological substances \\
\texttt{X65} & Intentional self-poisoning by and exposure to alcohol \\
\texttt{X66} & Intentional self-poisoning by and exposure to organic solvents and halogenated hydrocarbons and their vapours \\
\texttt{X67} & Intentional self-poisoning by and exposure to other gases and vapours \\
\texttt{X68} & Intentional self-poisoning by and exposure to pesticides \\
\texttt{X69} & Intentional self-poisoning by and exposure to other and unspecified chemicals and noxious substances \\
\texttt{X70} & Intentional self-harm by hanging, strangulation and suffocation \\
\texttt{X71} & Intentional self-harm by drowning and submersion \\
\texttt{X72} & Intentional self-harm by handgun discharge \\
\texttt{X73} & Intentional self-harm by rifle, shotgun and larger firearm discharge \\
\texttt{X74} & Intentional self-harm by other and unspecified firearm and gun discharge \\
\texttt{X75} & Intentional self-harm by explosive material \\
\texttt{X76} & Intentional self-harm by smoke, fire and flames \\
\texttt{X77} & Intentional self-harm by steam, hot vapours and hot objects \\
\texttt{X78} & Intentional self-harm by sharp object \\
\texttt{X79} & Intentional self-harm by blunt object \\
\texttt{X80} & Intentional self-harm by jumping from a high place \\
\texttt{X81} & Intentional self-harm by jumping or lying before moving object \\
\texttt{X82} & Intentional self-harm by crashing of motor vehicle \\
\texttt{X83} & Intentional self-harm by other specified means \\
\texttt{X84} & Intentional self-harm by unspecified means \\[6pt]

% ================================================================
% BLOCK X85-Y09  Assault
% ================================================================
\multicolumn{2}{l}{\textbf{X85--Y09 \quad Assault}} \\
\midrule
\texttt{X85} & Assault by drugs, medicaments and biological substances \\
\texttt{X86} & Assault by corrosive substance \\
\texttt{X87} & Assault by pesticides \\
\texttt{X88} & Assault by gases and vapours \\
\texttt{X89} & Assault by other specified chemicals and noxious substances \\
\texttt{X90} & Assault by unspecified chemical or noxious substance \\
\texttt{X91} & Assault by hanging, strangulation and suffocation \\
\texttt{X92} & Assault by drowning and submersion \\
\texttt{X93} & Assault by handgun discharge \\
\texttt{X94} & Assault by rifle, shotgun and larger firearm discharge \\
\texttt{X95} & Assault by other and unspecified firearm and gun discharge \\
\texttt{X96} & Assault by explosive material \\
\texttt{X97} & Assault by smoke, fire and flames \\
\texttt{X98} & Assault by steam, hot vapours and hot objects \\
\texttt{X99} & Assault by sharp object \\
\texttt{Y00} & Assault by blunt object \\
\texttt{Y01} & Assault by pushing from high place \\
\texttt{Y02} & Assault by pushing or placing victim before moving object \\
\texttt{Y03} & Assault by crashing of motor vehicle \\
\texttt{Y04} & Assault by bodily force \\
\texttt{Y05} & Sexual assault by bodily force \\
\texttt{Y06} & Neglect and abandonment \\
\texttt{Y07} & Other maltreatment syndromes \\
\texttt{Y08} & Assault by other specified means \\
\texttt{Y09} & Assault by unspecified means \\[6pt]

% ================================================================
% BLOCK Y10-Y34  Event of undetermined intent
% ================================================================
\multicolumn{2}{l}{\textbf{Y10--Y34 \quad Event of Undetermined Intent}} \\
\midrule
\texttt{Y10} & Poisoning by and exposure to nonopioid analgesics, antipyretics and antirheumatics, undetermined intent \\
\texttt{Y11} & Poisoning by and exposure to antiepileptic, sedative-hypnotic, antiparkinsonism and psychotropic drugs, not elsewhere classified, undetermined intent \\
\texttt{Y12} & Poisoning by and exposure to narcotics and psychodysleptics (hallucinogens), not elsewhere classified, undetermined intent \\
\texttt{Y13} & Poisoning by and exposure to other drugs acting on the autonomic nervous system, undetermined intent \\
\texttt{Y14} & Poisoning by and exposure to other and unspecified drugs, medicaments and biological substances, undetermined intent \\
\texttt{Y15} & Poisoning by and exposure to alcohol, undetermined intent \\
\texttt{Y16} & Poisoning by and exposure to organic solvents and halogenated hydrocarbons and their vapours, undetermined intent \\
\texttt{Y17} & Poisoning by and exposure to other gases and vapours, undetermined intent \\
\texttt{Y18} & Poisoning by and exposure to pesticides, undetermined intent \\
\texttt{Y19} & Poisoning by and exposure to other and unspecified chemicals and noxious substances, undetermined intent \\
\texttt{Y20} & Hanging, strangulation and suffocation, undetermined intent \\
\texttt{Y21} & Drowning and submersion, undetermined intent \\
\texttt{Y22} & Handgun discharge, undetermined intent \\
\texttt{Y23} & Rifle, shotgun and larger firearm discharge, undetermined intent \\
\texttt{Y24} & Other and unspecified firearm discharge, undetermined intent \\
\texttt{Y25} & Contact with explosive material, undetermined intent \\
\texttt{Y26} & Exposure to smoke, fire and flames, undetermined intent \\
\texttt{Y27} & Contact with steam, hot vapours and hot objects, undetermined intent \\
\texttt{Y28} & Contact with sharp object, undetermined intent \\
\texttt{Y29} & Contact with blunt object, undetermined intent \\
\texttt{Y30} & Falling, jumping or pushed from a high place, undetermined intent \\
\texttt{Y31} & Falling, lying or running before or into moving object, undetermined intent \\
\texttt{Y32} & Crashing of motor vehicle, undetermined intent \\
\texttt{Y33} & Other specified events, undetermined intent \\
\texttt{Y34} & Unspecified event, undetermined intent \\[6pt]

% ================================================================
% BLOCK Y35  Legal intervention (partial)
% ================================================================
\multicolumn{2}{l}{\textbf{Y35--Y36 \quad Legal Intervention and Operations of War (partial)}} \\
\midrule
\texttt{Y35} & Legal intervention \\*

\end{longtable}
\endgroup

\subsection*{A.2. Gold Reserves}

Based on \citet{sgb2024}, gold-bearing metallogenic provinces can be defined as regions where mineral occurrences and deposits appear in greater quantities or tonnage than neighboring areas, delimited by geologic-stratigraphic-tectonic boundaries, while districts are subdivisions of provinces, grouping spatially and genetically related deposits and occurrences. Therefore, we define \(Gold = 1\) for any municipality that overlaps at least one province or district of gold, considering that a district is not necessarily within a province. The polygons were obtained from the website of the Geological Survey of Brazil (SGB) in 2024.

\subsection*{A.3. Indigenous Territories}

The binary variable for indigenous territories was constructed based on the dataset of polygons of Indigenous territories made publicly available by National Foundation of Indigenous People (FUNAI) at GeoServer website. Therefore, we define \(IT = 1\) for any municipality that overlaps at least one indigenous territory in 2024.

\bibliographystyle{elsarticle-harv}
\bibliography{references}

\end{document}